\chapter{Introduction}

\section{Motivation}

Research activity is recorded across proceedings websites, bibliographic indexes,
author profiles, and institutional records. Each source describes only part of a
publication: an official page establishes a venue record and byline, while a
metadata index can add author identifiers, abstracts, references, and affiliation
strings. These records do not automatically agree, and an author name is not a
stable person identifier. As a result, counting papers alone can conceal the
collaboration structure that shapes research communities.

A scientific knowledge graph offers a way to represent those records and their
relationships explicitly. It can connect papers to researchers, organizations,
venues, and topics while retaining the source evidence and method that produced
each claim. This is useful for a research-program analyst who needs to ask not
only how much a venue publishes, but how collaboration groups connect, where
cross-sector co-publication occurs, and how the selected publication landscape
changes over time. Knowledge-graph design is relevant here because the graph's
meaning depends on explicit entity types, relation definitions, and provenance,
not merely on a collection of nodes and links \citep{hogan2022knowledgegraphs}.

AAAI and IJCAI provide a bounded, recognizable domain for this course project.
The selected records cover five consecutive years of main-track proceedings from
two established artificial-intelligence conference series. They provide both an
official publication boundary and enough repeated coauthorship to study
collaboration networks. The project does not treat these venues as a complete
sample of AI research. Instead, it uses the corpus as a defined window through
which to explore publication, collaboration, and citation relationships.

Network analysis can reveal patterns that publication totals do not show. A
researcher may appear in a large collaboration component, bridge communities, or
have a high number of distinct collaborators. An organization may co-publish with
partners in several sectors. A directed citation projection can identify papers
that receive many in-scope references. Such measures are structural summaries of
the selected graph; they do not by themselves demonstrate research quality,
causal influence, or a current business relationship. Scientific collaboration
networks are a well-established setting for this kind of structural analysis
\citep{newman2001collaboration}.

\section{Research questions}

The analysis is organized around four answerable questions tied to the available
release artifacts:

\begin{enumerate}[label=\textbf{RQ\arabic*.},leftmargin=3em,labelwidth=2.4em,labelsep=0.6em,align=left]
  \item How did the number and inferred topic mix of AAAI and IJCAI main-track
  papers in the selected corpus vary from 2021 through 2025?
  \item What collaboration structure is visible in the accepted researcher and
  organization projections, including components, coauthor degree, communities,
  and cross-sector publication support?
  \item Which researchers and papers rank highly under the project's configured
  betweenness and within-corpus citation measures, and how sensitive are these
  rankings to identity and evidence policies?
  \item What do the measured coverage, structural checks, AI-assessed edge
  reviews, historical topic evaluation, and identity-policy comparisons permit
  us to conclude about the reliability of those patterns?
\end{enumerate}

\section{Objectives and contributions}

The project has five practical objectives. First, it defines a reproducible paper
corpus from official venue records. Second, it links those records to OpenAlex
works through roster-derived DOI requests and explicit title/byline checks,
with historical DOI exceptions disclosed in \cref{chap:method}. Third, it builds a typed graph whose observed,
inferred-topic, curated-taxonomy, and derived relationships remain distinguishable.
Fourth, it computes named network projections and packages the measured results
as figures, tables, and source-linked findings. Fifth, it evaluates the artifact
while reporting limitations in identity resolution, affiliation coverage, topic
labels, and AI-assisted assessment.

The project's contribution is an auditable construction and analysis workflow,
not a new entity-resolution algorithm or a claim that the graph is a complete
record of AI research. Its main deliverables are a frozen v4 graph, versioned
analysis artifacts, reproducible static visualizations, structural validation
and database receipts, a read-only interactive Explorer, and a written account
that keeps candidate identity experiments separate from production results.

\section{Scope}

The production release contains \num{13617} roster-v10 papers from AAAI and
IJCAI main-track proceedings dated 2021--2025. The broader project source registry
is larger, but the release boundary is deliberately restricted to these two
series; unresolved or unavailable records from other eligible editions are
coverage limits rather than hidden graph members. The graph's five entity labels
are Researcher, Paper, Organization, Topic, and Venue. Its ten predicates are
defined in \cref{tab:relationship-schema}. The release graph is
\GraphVersion{}; experimental v5 identity candidates are discussed only as
sensitivity scenarios.

The report makes no inference about current employment, nationality, commercial
readiness, contract value, or causal impact. Organization sectors and countries
are metadata properties. Topic trends are based on inferred labels and exclude
papers for which the classifier abstained. Citation counts cover only references
whose target is another paper inside this corpus. These limits define how each
finding should be read.

\section{Report organization}

\Cref{chap:background} reviews knowledge graphs, scholarly networks, entity
resolution, topic classification, and graph-quality practice. \Cref{chap:method}
describes the data sources, matching and resolution rules, topic classifier, and
graph semantics. \Cref{chap:implementation} documents the software and Explorer
architecture. \Cref{chap:evaluation} defines the evaluation samples and
distinguishes structural validation from semantic assessment. \Cref{chap:results}
presents the graph, landscape, collaboration, centrality, citation, and sensitivity
results. The final chapters discuss interpretation and limitations, then conclude
with the supported findings and remaining submission details.
